% =========================
% Cas pratique 5
% =========================

\section[Cas 5]{Cas pratique 5 : La copropriété du brevet}

%--------------------------
\begin{frame}{Question 1}

\begin{block}{}
L'article \textbf{L613-29}b) CPI dispose que 
chacun des copropriétaires peut agir en contrefaçon à son seul profit. 
Le copropriétaire qui agit en contrefaçon doit notifier l'assignation délivrée
 aux autres copropriétaires ; il est sursis à statuer sur 
l'action tant qu'il n'est pas justifié de cette notification.\\
+ \textbf{L615-2} (titulaire apte à agir en cf. )
+ \textbf{L613-32} (pas de stipulation contraire)
+ \textbf{L614-9} (traduction FR de demande EP).
\end{block}

\begin{itemize}
  \item En l'espèce, la société DUVIN est copropriétaire du brevet EP avec Mme RAISIN.
  La société DUVIN est apte à agir en cf. Il n'y pas de stipulation contractuelles
  contraires puisque pas de convention de copropriété.
  \item En conclusion, la société DUVIN a le droit d'agir seule
  devant le tribunal judiciaire de Paris, à condition de notifier Mme. RAISIN.
\end{itemize}

\end{frame}

%--------------------------
\begin{frame}{Question 2}

\begin{block}{}
L'article \textbf{L615-7} CPI dispose que les dommages-intérêts sont fixés par:\\
a) manque à gagner, perte subies;\\
b) préjudice moral;\\
c) bénéfices du cf.\\
+ \textbf{L615-3} CPI (tout moyens)
+ \textbf{L615-5-2} CPI (droit de l'information).
\end{block}

\begin{itemize}
  \item En conclusion, la société DUVIN a à sa disposition
  pour jusitifier des D\&I du droit de l'information.
\end{itemize}

\end{frame}

\begin{frame}{Question 2}

\begin{block}{}
L'article \textbf{1240} du Code Civil dispose que
tout fait quelconque de l'homme, qui cause à autrui un dommage,
oblige celui par la faute duquel il est arrivé, à le réparer.
\end{block}

\begin{itemize}
  \item En conclusion, la société Duvin ne peut pas demander
  la réparation du préjudice qui est propre à Mme. Raisin.
  \item Si Mme. Raisin souhaite obtenir réparation,
  Mme. Raison peut agir en réparation de son préjudice propre
  par \textbf{L613-29}b) CPI.
\end{itemize}

\end{frame}

%--------------------------
\begin{frame}{Question 3}

\begin{block}{}
L'article \textbf{L614-12} CPI renvoie à l'article 138 CBE.\\
+ article \textbf{70} CPC (lien suffisant)\\
+ \textbf{CA Paris, pôle 5 ch. 2, 13 avril 2012, n°11/01329}\\
+ \textbf{CA Lyon, 12 sept. 2013, n°11/06334}.\\
\end{block}

\begin{itemize}
    \item Par hypothèse, la société Aubonpain est assignée.
    \item En conclusion, la société Aubonpain a le droit de
    demander une demande reconventionnelle en nullité à condition
    de la diriger contre l'ensemble des titulaires du brevet,
    à savoir Mme. Raisin et la société Duvin. A defaut, la
    demande reconventionnelle est nulle.
\end{itemize}

\end{frame}

%--------------------------
\begin{frame}{Question 4}

\begin{block}{}
L'article \textbf{L613-29}c) CPI dispose que chacun des copropriétaires 
peut concéder à un tiers une licence d'exploitation non exclusive à son profit,
 sauf à indemniser équitablement les autres copropriétaires 
 qui n'exploitent pas personnellement l'invention ou 
 qui n'ont pas concédé de licence d'exploitation.
  A défaut d'accord amiable,
cette indemnité est fixée par le tribunal judiciaire.
\end{block}

\begin{block}{}
L'article \textbf{L613-29}d) CPI dispose que'une 
licence d'exploitation exclusive 
ne peut être accordée qu'
avec l'accord de tous les copropriétaires ou par autorisation de justice. 
\end{block}

\begin{itemize}
    \item En conclusion, la société Duvin a le droit de concéder à 
    la société Aubonpain
     une licence d'exploitation non exclusive à condition d'indemniser
    Mme. Raisin. Toutefois, la société Duvin n'a pas le droit de concéder
    à la société Aubonpain une licence exclusive sans l'accord de Mme. Raisin.
\end{itemize}

\end{frame}